%!TEX TS-program = pdflatex
\documentclass[aspectratio=169,10pt]{beamer}
\usepackage[T2A]{fontenc}
\usepackage[utf8]{inputenc}
\usepackage[english,russian]{babel}
\usepackage{cmap}
\usepackage{amsmath,amssymb,mathtools,booktabs,array,etoolbox,PTSans,graphicx}
\graphicspath{{figures/}}
\renewcommand{\familydefault}{\sfdefault}
\renewcommand{\rmdefault}{\sfdefault}
\newbool{russian}\booltrue{russian}
\usepackage{HSE-theme/beamerthemeHSE}
\hypersetup{bookmarks=false}
\setbeamertemplate{blocks}[rounded][shadow=false]
\setbeamerfont{frametitle}{size=\Large}
\setbeamerfont{framesubtitle}{size=\small}
\setbeamerfont{block title}{size=\normalsize}
\setbeamerfont{block body}{size=\small}
\setlength{\abovedisplayskip}{5pt}
\setlength{\belowdisplayskip}{5pt}
\newcommand{\R}{\mathbb{R}}
\newcommand{\norm}[1]{\left\lVert #1\right\rVert}
\newcommand{\figwide}[1]{\begin{center}\includegraphics[height=.49\textheight,keepaspectratio]{#1}\end{center}}
\newcommand{\figshort}[1]{\begin{center}\includegraphics[height=.39\textheight,keepaspectratio]{#1}\end{center}}
\title[Методы оптимизации]{Методы оптимизации\\в машинном обучении}
\subtitle{Лекция 9. Метод Ньютона: от градиента к кривизне}
\author[Н. Лукьяненко]{Никита Лукьяненко}
\institute{Совместная программа СМОЛ ГУ и ФКН НИУ ВШЭ}
\date{2026/27 учебный год}
\begin{document}

% Блок 1: постановка задачи, 7 минут. Слайды 1--7.
% 1
\frame[plain]{\titlepage}

% 2
\begin{frame}{Где мы остановились}
\small
В предыдущих лекциях мы выбирали направление $-\nabla f(x)$ и отдельно подбирали длину шага $\alpha$.
\[
x_{k+1}=x_k-\alpha_k\nabla f(x_k).
\]
\begin{block}{Сегодняшний вопрос}
Можно ли использовать \textbf{кривизну} функции, чтобы сразу определить подходящее перемещение по каждой координате?
\end{block}
\end{frame}

% 3
\begin{frame}{Линейная модель видит наклон, но не кривизну}
\small
Разложение первого порядка даёт
\[
f(x+h)=f(x)+\nabla f(x)^\top h+o(\norm h).
\]
Его линейная часть не ограничена снизу по $h$, когда $\nabla f(x)\ne0$.
\begin{center}\includegraphics[width=.82\linewidth]{local_models.pdf}\end{center}
\end{frame}

% 4
\begin{frame}{Цели лекции}
\small
После лекции мы сможем:
\begin{enumerate}
\item вывести шаг Ньютона из квадратичной модели и линеаризации градиента;
\item доказать точное решение положительно определённой квадратичной задачи за один шаг;
\item понять, когда направление Ньютона убывает и когда полный шаг опасен;
\item доказать локальную квадратичную сходимость с явными предпосылками;
\item применить демпфирование и условие Армихо.
\end{enumerate}
\end{frame}

% 5
\begin{frame}{Градиент и гессиан: напоминание}
\small
Для дважды дифференцируемой $f:\R^n\to\R$:
\[
g(x)=\nabla f(x),\qquad H(x)=\nabla^2 f(x).
\]
Градиент задаёт изменение первого порядка, гессиан описывает изменение самого градиента.
\[
g(x+h)=g(x)+H(x)h+o(\norm h).
\]
\begin{alertblock}{Важное различие}
Гессиан может быть положительно определённым, неопределённым или вырожденным. От этого зависит смысл шага Ньютона.
\end{alertblock}
\end{frame}

% 6
\begin{frame}{Кривизна зависит от направления}
\small
Для единичного вектора $v$ величина $v^\top H(x)v$ задаёт вторую производную вдоль прямой $x+tv$.
\[
\left.\frac{d^2}{dt^2}f(x+tv)\right|_{t=0}=v^\top H(x)v.
\]
Если $H(x)\succ0$, кривизна положительна во всех ненулевых направлениях. Если $H(x)$ неопределён, локальная поверхность имеет направления разной кривизны.
\end{frame}

% 7
\begin{frame}{Маршрут на 80 минут}
\small
\begin{tabular}{@{}p{.70\linewidth}r@{}}
Постановка и связь с градиентным спуском & 7 мин\\
Квадратичная модель и вывод шага & 14 мин\\
Квадратичная функция: один шаг & 11 мин\\
Направление убывания и ограничения метода & 10 мин\\
Локальная квадратичная сходимость & 15 мин\\
Демпфирование и поиск Армихо & 15 мин\\
Примеры, сравнение и выводы & 8 мин
\end{tabular}
\end{frame}

% Блок 2: вывод шага, 14 минут. Слайды 8--17.
% 8
\begin{frame}{Когда у модели есть единственный минимум}
\small
Мы уже знаем из лекции 2 квадратичную модель функции в точке $x$:
\[
m_x(h)=f(x)+g(x)^\top h+\tfrac12 h^\top H(x)h.
\]
Здесь $h$ --- \textbf{перемещение} от текущей точки; $m_x(h)$ приближает значение $f(x+h)$.
Если $H(x)\succ0$, модель строго выпукла как функция перемещения $h$:
\[
\nabla_h m_x(h)=g(x)+H(x)h,
\qquad \nabla_h^2 m_x(h)=H(x)\succ0.
\]
Следующий вопрос: какое именно $h$ минимизирует эту модель?
\end{frame}

% 9
\begin{frame}{Стационарная точка модели --- перемещение, а не значение}
\small
Назовём $h_*$ \textbf{стационарной точкой модели}, если производная модели по перемещению равна нулю:
\[
\nabla_h m_x(h_*)=g(x)+H(x)h_*=0.
\]
При $H(x)\succ0$ это уравнение имеет ровно одно решение. Более того,
\[
m_x(h)-m_x(h_*)=\tfrac12(h-h_*)^\top H(x)(h-h_*)>0
\quad\text{для }h\ne h_*.
\]
Значит, \textbf{перемещение} $h_*$ даёт наименьшее значение \textbf{квадратичной модели} среди всех $h$. Для исходной функции $f(x+h)$ такого вывода пока нет.
\end{frame}

% 10
\begin{frame}{Вывод шага Ньютона из минимума модели}
\small
При $H(x)\succ0$ решаем условие стационарности модели:
\[
g(x)+H(x)h=0.
\]
Так как $H(x)$ обратима,
\[
\boxed{p_N(x)=-H(x)^{-1}g(x).}
\]
Вычислительно решают систему $H(x)p_N=-g(x)$, не формируя обратную матрицу.
\[
\boxed{x_{k+1}=x_k+p_N(x_k)}
\]
называется \textbf{полным шагом Ньютона}.
\end{frame}

% 11
\begin{frame}{Тот же шаг из уравнения стационарности}
\small
Искомый минимум удовлетворяет $g(x^*)=0$. Заменим градиент поблизости от $x$ линейной моделью:
\[
g(x+h)\approx g(x)+H(x)h.
\]
Принудительно обнулим приближение:
\[
g(x)+H(x)h=0
\quad\Longrightarrow\quad h=p_N(x).
\]
\begin{block}{Два взгляда на формулу}
При $H(x)\succ0$ шаг минимизирует квадратичную модель. При просто обратимом $H(x)$ он остаётся шагом Ньютона для уравнения $g(x)=0$, но может не вести к минимуму.
\end{block}
\end{frame}

% 12
\begin{frame}{Итерация в нескольких строках}
\small
В точке $x_k$:
\begin{enumerate}
\item вычислить $g_k=\nabla f(x_k)$ и $H_k=\nabla^2 f(x_k)$;
\item решить линейную систему $H_kp_k=-g_k$;
\item положить $x_{k+1}=x_k+p_k$.
\end{enumerate}
\begin{alertblock}{Условия нужны на каждом шаге}
Обратимость $H_k$ позволяет вычислить шаг. Положительная определённость дополнительно делает $p_k$ направлением убывания при $g_k\ne0$.
\end{alertblock}
\end{frame}

% 13
\begin{frame}{Одномерная формула}
\small
Для $f:\R\to\R$, если $f''(x)\ne0$,
\[
\boxed{x_{k+1}=x_k-\frac{f'(x_k)}{f''(x_k)}.}
\]
Например, для $f(x)=\frac12 a(x-x^*)^2$, $a>0$:
\[
f'(x)=a(x-x^*),\qquad f''(x)=a,
\]
поэтому $x_{k+1}=x^*$ независимо от начальной точки.
\end{frame}

% 14
\begin{frame}{Как кривизна меняет длину шага}
\small
В одном измерении $p_N=-f'/f''$. При одинаковом наклоне большой $f''>0$ даёт короткий шаг, маленький $f''>0$ даёт длинный.
\[
\underbrace{-\alpha f'(x)}_{\text{градиентный шаг}}
\qquad\text{и}\qquad
\underbrace{-\frac{f'(x)}{f''(x)}}_{\text{шаг Ньютона}}.
\]
В нескольких измерениях $H^{-1}$ корректирует разные направления по-разному.
\end{frame}

% 15
\begin{frame}{Геометрия двух направлений}
\small
\figwide{directions_contours.pdf}
Из синей начальной точки красная стрелка идёт по антиградиенту: она перпендикулярна линии уровня и показывает самое быстрое локальное убывание при малой длине шага. Зелёная стрелка --- перемещение Ньютона: оно учитывает вытянутость эллипсов и ведёт к их центру. Здесь функция квадратична, поэтому модель точна и зелёная стрелка попадает в минимум сразу.
\end{frame}

% 16
\begin{frame}{Гессиан как локальный масштаб координат}
\small
Если $H=Q\Lambda Q^\top\succ0$ и $g=\sum_i\gamma_iq_i$, то
\[
p_N=-Q\Lambda^{-1}Q^\top g
=-\sum_i\frac{\gamma_i}{\lambda_i}q_i.
\]
Здесь $q_i$ --- собственное направление гессиана, $\lambda_i>0$ --- кривизна в нём, $\gamma_i$ --- компонента градиента. Метод делит эту компоненту на кривизну: в крутом направлении перемещение короче, в пологом --- длиннее.

Если кривизна очень мала, полный шаг может оказаться слишком длинным за пределами локальной модели. Позже ограничим его длину поиском Армихо.
\end{frame}

% 17
\begin{frame}{Что пока доказано, а что ещё нет}
\small
Мы вывели формулу $p_N=-H^{-1}g$ и условия, при которых она минимизирует \textbf{локальную модель}.
\begin{itemize}
\item Убывает ли истинная функция при полном шаге?
\item Можно ли гарантировать сходимость из любой начальной точки?
\item Какова скорость вблизи невырожденного минимума?
\end{itemize}
Ответы различаются. Начнём с класса задач, где модель точна всюду.
\end{frame}

% Блок 3: квадратичная функция, 11 минут. Слайды 18--25.
% 18
\begin{frame}{Положительно определённая квадратичная задача}
\small
Пусть
\[
f(x)=\frac12x^\top Ax-b^\top x+c,
\qquad A=A^\top\succ0.
\]
Тогда
\[
g(x)=Ax-b,\qquad H(x)=A.
\]
Гессиан постоянен, поэтому разложение Тейлора второго порядка совпадает с $f(x+h)$ \textbf{точно}, без остатка.
\end{frame}

% 19
\begin{frame}{Единственный минимум квадратичной функции}
\small
Условие стационарности:
\[
Ax^*-b=0
\quad\Longrightarrow\quad
\boxed{x^*=A^{-1}b.}
\]
Поскольку $A\succ0$, этот стационарный вектор является единственным глобальным минимумом.
\[
f(x)-f(x^*)=\frac12(x-x^*)^\top A(x-x^*)>0
\quad\text{при }x\ne x^*.
\]
\end{frame}

% 20
\begin{frame}{Теорема: один шаг Ньютона}
\small
\begin{block}{Утверждение}
Для $f(x)=\frac12x^\top Ax-b^\top x+c$, $A=A^\top\succ0$, полный шаг Ньютона из любого $x$ приводит в единственный минимум $x^*=A^{-1}b$.
\end{block}
Действительно, направление решает систему
\[
Ap_N=-(Ax-b).
\]
Значит,
\[
p_N=A^{-1}b-x.
\]
\end{frame}

% 21
\begin{frame}{Доказательство одного шага: завершение}
\small
Подставляем полученное направление в полную итерацию:
\[
x_+=x+p_N=x+(A^{-1}b-x)=A^{-1}b=x^*.
\]
\[
\boxed{x_+=x^*.}\qquad\square
\]
\begin{alertblock}{Область действия}
Утверждение относится к точному решению линейной системы и постоянному положительно определённому гессиану. Численные ошибки и приближённое решение системы меняют буквальную точность одного шага.
\end{alertblock}
\end{frame}

% 22
\begin{frame}{Один шаг на карте линий уровня}
\small
\figwide{quadratic_one_step.pdf}
Для квадратичной функции локальная модель совпадает с исходной задачей. Поэтому полный шаг Ньютона сразу попадает в центр эллипсов.
\end{frame}

% 23
\begin{frame}{Численный пример в двух координатах}
\small
Возьмём
\[
A=\begin{pmatrix}4&0\\0&1\end{pmatrix},\quad
b=\begin{pmatrix}4\\-2\end{pmatrix},\quad
x_0=\begin{pmatrix}3\\2\end{pmatrix}.
\]
Тогда $x^*=(1,-2)^\top$ и $g_0=Ax_0-b=(8,4)^\top$.
\[
p_0=-A^{-1}g_0=\begin{pmatrix}-2\\-4\end{pmatrix},
\qquad x_1=x_0+p_0=\begin{pmatrix}1\\-2\end{pmatrix}.
\]
Проверка: $g(x_1)=0$.
\end{frame}

% 24
\begin{frame}{Почему градиентному спуску нужен выбор шага}
\small
Для той же квадратичной функции градиентный шаг имеет вид
\[
x_+=x-\alpha(Ax-b).
\]
Чтобы он попал в $x^*$ за один шаг, необходимо одновременно выполнить $\alpha=1/\lambda_i$ во всех активных собственных направлениях.
\begin{block}{Вывод}
При разных собственных значениях одно число $\alpha$ обычно не может устранить всю ошибку. Матрица $A^{-1}$ корректирует каждую компоненту отдельно.
\end{block}
\end{frame}

% 25
\begin{frame}{Зависимость результата от структуры задачи}
\small
\begin{tabular}{@{}p{.30\linewidth}p{.59\linewidth}@{}}
\toprule
Функция & Полный шаг Ньютона\\
\midrule
$A\succ0$ и $H(x)=A$ постоянно & точное решение за один шаг\\
Общая $C^2$ функция около минимума & квадратичная локальная сходимость при дополнительных условиях\\
Далеко от минимума & полный шаг может увеличить $f$\\
\bottomrule
\end{tabular}

\par\vspace{.7em}
Далее разберём вторую и третью строки строго.
\end{frame}

% Блок 4: убывание и ограничения, 10 минут. Слайды 26--32.
% 26
\begin{frame}{Когда направление Ньютона убывает}
\small
В лекции 8 направление $p$ называлось направлением убывания, когда $g(x)^\top p<0$: при малом положительном $\alpha$ функция $f(x+\alpha p)$ становится меньше.
\begin{block}{Утверждение для метода Ньютона}
Если $H(x)\succ0$ и $g(x)\ne0$, то $p_N=-H(x)^{-1}g(x)$ является строгим направлением убывания:
\[
\boxed{g(x)^\top p_N<0.}
\]
\end{block}
Это утверждение касается \textbf{направления}. Полный шаг $\alpha=1$ может увеличить функцию; такой пример увидим далее.
\end{frame}

% 27
\begin{frame}{Доказательство направления убывания}
\small
Положительная определённость $H$ влечёт $H^{-1}\succ0$. Поэтому для $g\ne0$
\[
g^\top p_N=-g^\top H^{-1}g<0.
\]
Если $\varphi(\alpha)=f(x+\alpha p_N)$, то
\[
\varphi'(0)=g^\top p_N<0.
\]
По определению производной найдётся $\bar\alpha>0$ такое, что $f(x+\alpha p_N)<f(x)$ при $0<\alpha\le\bar\alpha$. $\square$
\end{frame}

% 28
\begin{frame}{Что происходит при нулевом градиенте}
\small
Если $g(x)=0$ и $H(x)$ обратима, то $p_N=0$: метод уже находится в стационарной точке.
\begin{itemize}
\item При $H(x)\succ0$ это строгий локальный минимум.
\item При $H(x)\prec0$ это строгий локальный максимум.
\item При неопределённом $H(x)$ это седловая точка.
\end{itemize}
Формула шага сама по себе не классифицирует стационарную точку.
\end{frame}

% 29
\begin{frame}{Неопределённый гессиан: шаг к седлу}
\small
Рассмотрим $f(x,y)=\frac12(x^2-y^2)$ в точке $(0,1)$.
\[
g=(0,-1)^\top,\quad H=\operatorname{diag}(1,-1),\quad p_N=(0,-1)^\top.
\]
Тогда $g^\top p_N=1>0$, то есть это \textbf{направление роста}.
\[
f(0,1)=-\frac12,\qquad f(0,0)=0.
\]
Полный шаг приходит в седловую точку и увеличивает функцию.
\end{frame}

% 30
\begin{frame}{Вырожденный гессиан: шаг не определён}
\small
Для $f(x)=x^3-3x$ в точке $x=0$:
\[
f'(0)=-3,\qquad f''(0)=0.
\]
Уравнение Ньютона $f''(0)p=-f'(0)$ превращается в $0\cdot p=3$ и не имеет решения.
\begin{block}{Практическое следствие}
Перед вычислением направления проверяют возможность решения системы и свойства гессиана. При проблемах используют модификацию гессиана или другое направление.
\end{block}
\end{frame}

% 31
\begin{frame}{Полный шаг может провалиться даже при $H>0$}
\small
Пусть $f(x)=\sqrt{1+x^2}$. Здесь
\[
f'(x)=\frac{x}{\sqrt{1+x^2}},\qquad
f''(x)=\frac{1}{(1+x^2)^{3/2}}>0.
\]
Шаг Ньютона даёт $x_+=-x^3$. Из $x_0=2$ получаем $x_1=-8$ и $f(x_1)=\sqrt{65}>\sqrt5=f(x_0)$.
\begin{center}\includegraphics[width=.83\linewidth]{full_step_failure.pdf}\end{center}
\end{frame}

% 32
\begin{frame}{Почему локальное направление не гарантирует полный шаг}
\small
Если $g^\top p_N<0$, производная $\varphi'(0)$ отрицательна. Это гарантирует убывание только при \textbf{достаточно малой} длине шага.
\[
\varphi(\alpha)=f(x+\alpha p_N),\qquad \varphi'(0)<0.
\]
При $\alpha=1$ мы можем выйти за область, где квадратичная модель описывает $f$ достаточно точно. Нужен механизм выбора $\alpha$.
\end{frame}

% Блок 5: локальная теорема, 15 минут. Слайды 33--41.
% 33
\begin{frame}{Теорема о локальной квадратичной сходимости}
\footnotesize
\begin{block}{Теорема}
Пусть $f\in C^2$ в окрестности $x^*$, $\nabla f(x^*)=0$, а $H^*=\nabla^2f(x^*)\succ0$.
Предположим, что в некотором шаре $B_r(x^*)$ гессиан липшицев и ограничен снизу:
\[
H(x)\succeq mI\quad(m>0),\qquad
\norm{H(x)-H(y)}_2\le L_H\norm{x-y}_2.
\]
Тогда для достаточно близкого $x_0$ полный метод Ньютона корректен, остаётся в этом шаре и
\[
\boxed{\norm{x_{k+1}-x^*}_2\le\frac{L_H}{2m}\norm{x_k-x^*}_2^2.}
\]
\end{block}
Если $L_H=0$, гессиан постоянен в шаре и следующий шаг сразу даёт $x^*$.
\end{frame}

% 34
\begin{frame}{Доказательство: почему гессиан обратим поблизости}
\small
\begin{block}{Доказательство, шаг 1 из 5}
Предпосылка $H^*\succ0$ означает $\lambda_{\min}(H^*)>0$.
Поскольку $H(x)$ непрерывен, можно уменьшить шар около $x^*$ так, чтобы
\[
\norm{H(x)-H^*}_2\le\tfrac12\lambda_{\min}(H^*)
\]
и, следовательно,
\[
H(x)\succeq \tfrac12\lambda_{\min}(H^*)I.
\]
\end{block}
Липшицевость гессиана нужна отдельно для \textbf{квадратичной оценки ошибки}; одной непрерывности для этой оценки недостаточно.
\end{frame}

% 35
\begin{frame}{Доказательство: точная формула для градиента}
\small
\textbf{Шаг 2 из 5.} Обозначим $e=x-x^*$. По основной теореме анализа вдоль отрезка:
\[
g(x)-g(x^*)=\int_0^1H(x^*+te)e\,dt.
\]
Так как $g(x^*)=0$,
\[
\boxed{g(x)=\int_0^1H(x^*+te)e\,dt.}
\]
Эта формула точна. Именно она позволяет контролировать ошибку после шага без неформального символа $\approx$.
\end{frame}

% 36
\begin{frame}{Доказательство: ошибка после шага}
\small
\textbf{Шаг 3 из 5.} Из $x_+=x-H(x)^{-1}g(x)$ и $e=x-x^*$ получаем
\[
e_+=e-H(x)^{-1}g(x)
=H(x)^{-1}\bigl(H(x)e-g(x)\bigr).
\]
Подставляя интегральную формулу:
\[
\boxed{e_+=H(x)^{-1}\int_0^1
\bigl[H(x)-H(x^*+te)\bigr]e\,dt.}
\]
Для квадратичной функции $H$ постоянно, интеграл равен нулю.
\end{frame}

% 37
\begin{frame}{Доказательство: почему ошибка возводится в квадрат}
\footnotesize
\textbf{Шаг 4 из 5.} Из $H(x)\succeq mI$ следует $\norm{H(x)^{-1}}_2\le 1/m$. Липшицевость даёт
\[
\norm{H(x)-H(x^*+te)}_2
\le L_H\norm{x-(x^*+te)}_2=L_H(1-t)\norm e_2.
\]
Поэтому
\[
\begin{aligned}
\norm{e_+}_2
&\le\frac1m\int_0^1 L_H(1-t)\norm e_2^2\,dt\\
&=\boxed{\frac{L_H}{2m}\norm e_2^2.}
\end{aligned}
\]
Это центральная оценка теоремы. Осталось проверить, что итерации остаются в выбранной окрестности.
\end{frame}

% 38
\begin{frame}{Доказательство: итерации остаются рядом}
\small
\textbf{Шаг 5 из 5.} Пусть $L_H>0$ и выберем
\[
0<\rho<r,\qquad \rho\le m/L_H.
\]
Если $\norm{e_0}_2\le\rho$, то из оценки следует
\[
\norm{e_{k+1}}_2\le\frac{L_H}{2m}\rho\norm{e_k}_2
\le\tfrac12\norm{e_k}_2.
\]
Индукция показывает, что все точки остаются в шаре $B_\rho(x^*)$ и сходятся к $x^*$. Внутри него действует более сильная квадратичная оценка. $\square$
\end{frame}

% 39
\begin{frame}{Что означает квадратичная сходимость}
\small
Если $e_k=\norm{x_k-x^*}$, то рядом с решением
\[
e_{k+1}\le C e_k^2,\qquad C=\frac{L_H}{2m}.
\]
\textbf{Пример:} $f(x)=e^x-x-1$, минимум $x^*=0$. Ньютон даёт $x_{k+1}=x_k-1+e^{-x_k}$ при $x_0=0{,}5$.
\begin{center}
\begin{tabular}{@{}rcc@{}}
\toprule
$k$ & $e_k=|x_k|$ & $e_{k+1}/e_k^2$\\
\midrule
0 & $0{,}5$ & $0{,}426$\\
1 & $0{,}106531$ & $0{,}483$\\
2 & $0{,}005478$ & $0{,}499$\\
3 & $1{,}50\cdot10^{-5}$ & $0{,}500$\\
4 & $1{,}12\cdot10^{-10}$ & ---\\
\bottomrule
\end{tabular}
\end{center}
Отношение стремится к $1/2$: новая ошибка примерно равна половине квадрата предыдущей.
\end{frame}

% 40
\begin{frame}{Невырожденность важна: пример $x^4$}
\small
Для $f(x)=x^4$ минимум находится в $x^*=0$, но $f''(0)=0$.
При $x\ne0$:
\[
p_N=-\frac{4x^3}{12x^2}=-\frac{x}{3},\qquad
x_+=\frac23x.
\]
Следовательно, $|x_k|=(2/3)^k|x_0|$: сходимость \textbf{линейная}, а не квадратичная.
\figshort{singular_linear_rate.pdf}
\end{frame}

% 41
\begin{frame}{Границы локальной теоремы}
\small
Теорема не утверждает глобальной сходимости из произвольной точки. Она требует:
\begin{itemize}
\item стационарный $x^*$ с $H(x^*)\succ0$;
\item липшицев гессиан рядом с $x^*$;
\item достаточно близкое начало для полного шага.
\end{itemize}
Вне этой окрестности полезно выбирать длину шага по фактическому значению $f$.
\end{frame}

% Блок 6: демпфирование и Армихо, 15 минут. Слайды 42--50.
% 42
\begin{frame}{Демпфированный метод Ньютона}
\small
\textbf{Демпфированный} означает «с уменьшенной длиной шага». Направление Ньютона сохраняется, но его умножают на коэффициент $0<\alpha_k\le1$:
\[
\boxed{x_{k+1}=x_k+\alpha_k p_k,\qquad H_kp_k=-g_k.}
\]
При $\alpha_k=1$ получаем полный шаг, при $\alpha_k<1$ --- демпфированный. Если $H_k\succ0$ и $g_k\ne0$, направление убывает. Теперь выбираем длину по фактическому уменьшению функции.
\end{frame}

% 43
\begin{frame}{Условие Армихо}
\small
Фиксируем $c\in(0,1)$ и $\beta\in(0,1)$. Пробуем $\alpha=1,\beta,\beta^2,\ldots$, пока
\[
\boxed{f(x+\alpha p)\le f(x)+c\alpha g^\top p.}
\]
Поскольку $g^\top p<0$, правая часть меньше $f(x)$. Условие требует измеримого убывания, пропорционального предсказанному наклону.
\end{frame}

% 44
\begin{frame}{Почему поиск Армихо завершается}
\footnotesize
Пусть $g^\top p<0$, а градиент имеет константу Липшица $L$ вдоль пробуемого отрезка. Лемма о спуске даёт
\[
f(x+\alpha p)\le f(x)+\alpha g^\top p+\frac{L\alpha^2}{2}\norm p^2.
\]
Условие Армихо выполнено, если
\[
0<\alpha\le\frac{2(1-c)(-g^\top p)}{L\norm p^2}.
\]
Геометрическая последовательность $1,\beta,\beta^2,\ldots$ в конце концов станет меньше этой положительной границы. $\square$
\end{frame}

% 45
\begin{frame}{Армихо исправляет слишком длинный шаг}
\small
Для $f(x)=\sqrt{1+x^2}$ в точке $x=2$ направление Ньютона равно $p=-10$.
При $c=0.1$, $\beta=1/2$ шаги $1$ и $1/2$ отклоняются, $\alpha=1/4$ принимается: $x_+=-1/2$.
\begin{center}\includegraphics[width=.83\linewidth]{armijo_example.pdf}\end{center}
\end{frame}

% 46
\begin{frame}{Алгоритм демпфированного Ньютона}
\footnotesize
\begin{block}{На каждой итерации}
\begin{enumerate}
\item Вычислить $g_k$ и $H_k$. Если $\norm{g_k}$ мала, остановиться.
\item Получить направление $p_k$ из $H_kp_k=-g_k$. Проверить $g_k^\top p_k<0$.
\item Начать с $\alpha=1$. Пока не выполнено условие Армихо, заменить $\alpha\leftarrow\beta\alpha$.
\item Обновить $x_{k+1}=x_k+\alpha p_k$.
\end{enumerate}
\end{block}
Если $H_k$ не положительно определён, перед поиском шага нужно получить другое направление убывания, например из положительно определённой модификации $H_k$.
\end{frame}

% 47
\begin{frame}{Пример траектории с демпфированием}
\small
\figwide{damped_progress.pdf}
В примере $f(x)=\sqrt{1+x^2}$ поиск Армихо удерживает итерации от резкого прыжка $2\mapsto-8$ и приводит к минимуму $x^*=0$.
\end{frame}

% 48
\begin{frame}{Почему вблизи минимума принимается $\alpha=1$}
\footnotesize
Пусть $H(x)\succeq mI$, а гессиан липшицев с константой $L_H$ на отрезке от $x$ до $x+p$. Тогда остаток формулы Тейлора можно \textbf{оценить явно}:
\[
f(x+p)=f(x)+g^\top p+\tfrac12p^\top H(x)p+R,
\qquad |R|\le\tfrac{L_H}{6}\norm p^3.
\]
Для шага Ньютона $g=-Hp$, поэтому $g^\top p=-p^\top Hp$. Условие Армихо при $\alpha=1$ выполнится, если
\[
R\le(\tfrac12-c)p^\top Hp.
\]
При $c<1/2$ правая часть не меньше $(\tfrac12-c)m\norm p^2$, а $|R|$ имеет порядок $\norm p^3$. Поэтому достаточно малый полный шаг принимается.
\end{frame}

% 49
\begin{frame}{Демпфирование сохраняет быстрый локальный режим}
\small
Если выполнены предпосылки локальной теоремы и $c\in(0,1/2)$, то в достаточно малой окрестности $x^*$ поиск Армихо принимает $\alpha=1$.
После этого метод совпадает с полным методом Ньютона:
\[
\norm{x_{k+1}-x^*}\le\frac{L_H}{2m}\norm{x_k-x^*}^2.
\]
Демпфирование контролирует шаг вдали от решения, а локальная квадратичная скорость возвращается рядом с ним.
\end{frame}

% 50
\begin{frame}{Цена одной итерации}
\small
На каждом шаге нужны градиент, гессиан и решение линейной системы. Для плотной матрицы размерности $n$ факторизация обычно требует порядка $n^3$ операций, хранение гессиана --- порядка $n^2$ чисел.
\begin{block}{Мост к следующей лекции}
Можно ли сохранить коррекцию разных масштабов задачи, но заменить дорогой гессиан более простой положительно определённой матрицей? Так мы придём к \textbf{предобусловливанию} градиентного спуска.
\end{block}
\end{frame}

% Блок 7: приложения и закрепление, 8 минут. Слайды 51--56.
% 51
\begin{frame}{Пример из ML: регуляризованная логистическая регрессия}
\footnotesize
В лекции 7 мы уже находили градиент логистической потери. Теперь посмотрим, какую дополнительную информацию даёт гессиан. Для $y_i\in\{-1,1\}$ и $a_i\in\R^n$:
\[
f(w)=\sum_{i=1}^N\log(1+e^{-y_i a_i^\top w})
+\frac{\lambda}{2}\norm w^2,\qquad \lambda>0.
\]
Положив $s_i=(1+e^{y_i a_i^\top w})^{-1}$, получаем
\[
H(w)=\sum_{i=1}^N s_i(1-s_i)a_i a_i^\top+\lambda I\succ0,
\]
поскольку $s_i(1-s_i)>0$ и $\lambda I\succ0$. Поэтому $H(w)p=-g(w)$ даёт направление убывания при $g(w)\ne0$; Армихо подбирает длину. Полное численное решение системы может быть дорого при большом числе признаков.
\end{frame}

% 52
\begin{frame}{Пример из ML: метод наименьших квадратов}
\small
Эту задачу мы изучали в лекции 3. Для $f(w)=\frac12\norm{Xw-y}_2^2$:
\[
g(w)=X^\top(Xw-y),\qquad H=X^\top X.
\]
Если столбцы $X$ линейно независимы, $H\succ0$ постоянно. Направление удовлетворяет
\[
X^\top Xp=-X^\top(Xw-y)
\quad\Longrightarrow\quad X^\top X(w+p)=X^\top y.
\]
Значит, один \textbf{полный} шаг даёт точный минимум в точной арифметике.
\begin{alertblock}{Численная оговорка}
На практике МНК часто решают через QR или SVD, не формируя $X^\top X$: так устойчивее при плохо обусловленной матрице $X$.
\end{alertblock}
\end{frame}

% 53
\begin{frame}{Как выбирать действие по состоянию гессиана}
\small
\begin{tabular}{@{}p{.29\linewidth}p{.60\linewidth}@{}}
\toprule
Состояние & Действие\\
\midrule
$H\succ0$, точка далеко от минимума & направление Ньютона и поиск Армихо\\
$H\succ0$, точка близко к минимуму & полный шаг; ожидаем квадратичную скорость\\
$H$ неопределён или вырожден & модифицировать гессиан либо выбрать другое направление убывания\\
\bottomrule
\end{tabular}

\par\vspace{.5em}
Критерий близости здесь теоретический; на практике его заменяет успешное принятие полного шага.
\end{frame}

% 54
\begin{frame}{Короткая проверка понимания}
\small
\begin{enumerate}
\item Почему полный метод решает положительно определённую квадратичную задачу за один шаг?
\item Какое свойство $H(x)$ делает $p_N$ направлением убывания?
\item Почему $f\in C^2$ и $H(x^*)\succ0$ сами по себе не дают выписанную квадратичную оценку?
\item При каком условии на $c$ Армихо в конце концов принимает полный шаг рядом с $x^*$?
\end{enumerate}
\end{frame}

% 55
\begin{frame}{Итоги}
\small
\[
\boxed{H_kp_k=-g_k,\qquad x_{k+1}=x_k+\alpha_kp_k.}
\]
\begin{itemize}
\item При $H\succ0$ направление Ньютона убывает; квадратичную функцию оно решает за один полный шаг.
\item При липшицевом гессиане и невырожденном минимуме полный метод сходится локально квадратично.
\item Поиск Армихо контролирует длину шага и рядом с минимумом допускает $\alpha=1$ при $c<1/2$.
\end{itemize}
\vspace{.4em}
\footnotesize Основной источник оформления: лекции 2--5 курса «Методы оптимизации в машинном обучении»; математический вывод выполнен в этой лекции.
\end{frame}

% 56
\begin{frame}{Лекция 10. Предобусловливание и геометрия методов оптимизации}
\footnotesize
\begin{block}{Почему эта тема следующая}
В лекциях 2--4 мы увидели, как форма линий уровня и спектр влияют на градиентный спуск; в лекции 6 оценили его скорость. Лекции 7--8 дали вычисление градиента и выбор шага. Сегодня метод Ньютона показал, как кривизна устраняет растяжение на квадратичной задаче, но точный гессиан может быть дорогим.
\end{block}
\textbf{Идея:} выбрать более простую матрицу $P\succ0$, отражающую масштаб задачи, и изменить направление:
\[
\boxed{x_{k+1}=x_k-\alpha_k P^{-1}\nabla f(x_k).}
\]
При $P=I$ это обычный градиентный спуск; если на каждом шаге брать $P=H(x_k)\succ0$ и $\alpha_k=1$, получаем шаг Ньютона.
\par\vspace{.4em}
\textbf{План лекции 10:}
\begin{enumerate}
\item Замена координат и новая геометрия квадратичной функции.
\item Матрица предобусловливания и предобусловленный градиентный спуск.
\item Как спектр после замены координат меняет скорость сходимости.
\item Простые варианты $P$ и сравнение с GD и методом Ньютона.
\end{enumerate}
\end{frame}

\end{document}
